\section{Cây lan ý và các thông số chăm sóc}
Cây lan ý (\textit{Spathiphyllum}), còn được biết đến với tên tiếng Anh là peace lily, là loài cây nội thất phổ biến nhờ hình thái lá xanh bóng, hoa trắng dạng mo và khả năng thích nghi tương đối tốt trong nhà. Tài liệu chăm sóc cây cho thấy lan ý có thể sống trong điều kiện ánh sáng yếu, song cây sinh trưởng và ra hoa tốt hơn khi nhận ánh sáng gián tiếp sáng; cây cũng ưa môi trường ấm, ẩm và cần tránh ánh nắng trực tiếp mạnh \cite{peace_lily_care_bhg}. Những đặc điểm này phù hợp với căn hộ chung cư, nơi cây thường được đặt ở gần cửa sổ, ban công có che chắn hoặc khu vực sinh hoạt có ánh sáng khuếch tán.

Về vai trò đối với IAQ, lan ý cần được nhìn nhận thận trọng. Cây có thể là thành phần hỗ trợ trong hệ kiểm soát môi trường, nhưng không thay thế thông gió, máy lọc không khí hoặc các thiết bị xử lý ô nhiễm chuyên dụng. Tài liệu về hệ thủy canh cho IAQ xếp lan ý vào nhóm cây có tiềm năng hỗ trợ hấp thụ CO2 và VOC; khi kết hợp với IoT, các thông số như CO2, nhiệt độ, độ ẩm, ánh sáng và dinh dưỡng có thể được theo dõi nhằm duy trì điều kiện sinh trưởng ổn định cho cây \cite{hydroponic_iaq_survey}. Vì vậy, khóa luận sử dụng IoT để giám sát lan ý như một phần của môi trường sống, đồng thời liên hệ dữ liệu chăm sóc cây với dữ liệu IAQ xung quanh.

\subsection{TDS và pH}
TDS biểu thị tổng lượng chất rắn hòa tan trong nước hoặc dung dịch dinh dưỡng. Trong thủy canh, EC thường được sử dụng như đại lượng quản lý dinh dưỡng vì các muối hòa tan tạo thành ion và làm thay đổi khả năng dẫn điện của dung dịch. EC quá cao có thể gây stress thẩm thấu, mất cân bằng dinh dưỡng hoặc độc ion; ngược lại, EC quá thấp thường phản ánh tình trạng dung dịch nghèo dinh dưỡng \cite{singh_hydroponics_ph_ec}. Trong nguyên mẫu của khóa luận, cảm biến TDS được dùng như một chỉ báo thực dụng để phát hiện dung dịch quá loãng hoặc quá đậm, từ đó gợi ý người dùng thay nước, bổ sung nước hoặc kiểm tra lại dung dịch dinh dưỡng.

pH phản ánh tính axit hoặc kiềm của dung dịch. Đối với lan ý thủy canh quy mô hộ gia đình, tài liệu hướng dẫn khuyến nghị duy trì dung dịch hơi axit, khoảng pH 6.0--6.5 \cite{peace_lily_hydro_brightlane}; trong nghiên cứu thủy canh trên \textit{Spathiphyllum}, dung dịch được duy trì ở pH 5.8 và EC 1.2 dS/m trong các thí nghiệm sinh trưởng \cite{peace_lily_ec_ppf}. Khi pH lệch khỏi vùng phù hợp, cây vẫn có nước và ánh sáng nhưng khả năng hấp thụ một số chất dinh dưỡng có thể suy giảm. Vì vậy, đo pH giúp hệ thống đánh giá trạng thái dung dịch và đưa ra cảnh báo chăm sóc cây kịp thời hơn.

\subsection{Ánh sáng}
Ánh sáng quyết định trực tiếp đến quang hợp và gián tiếp ảnh hưởng đến hình thái sinh trưởng của cây. Lan ý có thể tồn tại trong điều kiện ánh sáng thấp, nhưng ánh sáng gián tiếp sáng thường giúp cây sinh trưởng và ra hoa tốt hơn; ngược lại, nắng trực tiếp mạnh có thể làm hư hại lá \cite{peace_lily_care_bhg}. Trong nghiên cứu thủy canh trên \textit{Spathiphyllum}, mức PPF 70--100 $\mu$mol m$^{-2}$ s$^{-1}$ cho đáp ứng sinh trưởng và quang hợp tốt hơn so với các mức quá thấp hoặc quá cao \cite{peace_lily_ec_ppf}. Do vị trí đặt cây trong căn hộ có thể bị che khuất hoặc xa nguồn sáng tự nhiên, cảm biến ánh sáng được sử dụng để đánh giá điều kiện chiếu sáng và làm cơ sở điều khiển đèn bổ sung.

\subsection{Liên hệ giữa chăm sóc cây và IAQ}
Trong hệ thống đề xuất, dữ liệu chăm sóc cây và dữ liệu IAQ không được xử lý như hai nhóm thông tin tách rời. Ánh sáng thấp kéo dài có thể làm giảm quang hợp; nhiệt độ hoặc độ ẩm không phù hợp vừa ảnh hưởng đến cây, vừa ảnh hưởng đến cảm giác tiện nghi của người sử dụng. Khi CO2 tăng cao, hệ thống ưu tiên cảnh báo thông gió hoặc bật quạt; khi ánh sáng không đủ, hệ thống có thể đề xuất hoặc kích hoạt đèn bổ sung. Nhờ đó, các ngưỡng điều khiển được xây dựng để phục vụ đồng thời sức khỏe cây và chất lượng môi trường sống, thay vì chỉ tập trung vào nước và dinh dưỡng như nhiều hệ thống trồng cây thông minh đơn lẻ.
